% Source coverage: physical PDF pp. 315--317 (printed pp. 292--294), beginning at the Problems heading and continuing through the complete chapter reference list.
% Problems: 5. Unnumbered displays: 2. References: 1--22, with no lettered reference identifiers.
% Corrected errata: the spectral-function definitions include their intermediate states and momentum delta functions; Ref. 15 uses the paper's publication year, 1977.

\chapterbackmatter{Problems}

\begin{enumerate}
\item Consider a theory of a fermion field $\psi$ interacting with a scalar field $\phi$ with interactions of the form $\bar\psi\psi\phi$ and $\phi^4$. List the operators that appear in the operator product expansion of $\bar\psi(x)\psi(0)$ with coefficient functions that (judging from perturbation theory) are singular and non-vanishing for $x\to0$. Describe how you would calculate these coefficient functions to one-loop order.

\item Consider quantum chromodynamics with $N$ massless quarks, and define the spectral functions of the scalar and pseudoscalar quark bilinears by
\[
\begin{aligned}
&\sum_N\delta^4(p-p_N)
\bra{\Omega}\bar\psi(0)\lambda_\alpha\psi(0)\ket{N}
\bra{\Omega}\bar\psi(0)\lambda_\beta\psi(0)\ket{N}^{*}
\\
&\hspace{36mm}
=(2\pi)^{-3}\theta(p^0)\rho_{\alpha\beta}^{S}(-p^2),
\end{aligned}
\]
\[
\begin{aligned}
&\sum_N\delta^4(p-p_N)
\bra{\Omega}\bar\psi(0)\gamma_5\lambda_\alpha\psi(0)\ket{N}
\bra{\Omega}\bar\psi(0)\gamma_5\lambda_\beta\psi(0)\ket{N}^{*}
\\
&\hspace{36mm}
=(2\pi)^{-3}\theta(p^0)\rho_{\alpha\beta}^{P}(-p^2),
\end{aligned}
\]
where $\lambda_\alpha$ are a complete set of traceless Hermitian $N\times N$ matrices, normalized so that $\operatorname{Tr}(\lambda_\alpha\lambda_\beta)=2\delta_{\alpha\beta}$. What spectral function sum rules are satisfied by linear combinations of $\rho_{\alpha\beta}^{S}(\mu^2)$ and $\rho_{\alpha\beta}^{P}(\mu^2)$?

\item Derive Eq.~\eqref{eq:20.6.8} in the parton model from the formulas \hyperref[eq:8.7.7]{(8.7.7)} and \hyperref[eq:8.7.38]{(8.7.38)} for Compton scattering.

\item List the gauge-invariant symmetric traceless tensors of twist four in quantum chromodynamics.

\item In the massless scalar field theory with interaction $-g\phi^4/24$ (with $g>0$), where in the complex plane would you expect the function \eqref{eq:20.7.3} to have renormalon singularities?
\end{enumerate}

\chapterbackmatter{References}

\begin{enumerate}
\item \phantomsection\label{ch20-ref-1}
K. Wilson, \emph{Phys. Rev.} \textbf{179}, 1499 (1969).

\item \phantomsection\label{ch20-ref-2}
W. Zimmerman, in \emph{Lectures on Elementary Particles and Quantum Field Theory -- 1970 Brandeis University Summer Institute in Theoretical Physics}, eds. S. Deser, H. Pendleton, and M. Grisaru, (MIT Press, Cambridge, 1970).

\item \phantomsection\label{ch20-ref-3}
S. Weinberg, \emph{Phys. Rev.} \textbf{118}, 838 (1960).

\item \phantomsection\label{ch20-ref-4}
C. Bernard, A. Duncan, J. LoSecco, and S. Weinberg, \emph{Phys. Rev.} \textbf{D12}, 792 (1975).

\item \phantomsection\label{ch20-ref-5}
S. Weinberg, \emph{Phys. Rev. Lett.} \textbf{18}, 507 (1967). For the general case, see \hyperref[ch20-ref-4]{Ref.~4}.

\item \phantomsection\label{ch20-ref-6}
K. Kawarabayashi and M. Suzuki, \emph{Phys. Rev. Lett.} \textbf{16}, 255 (1966); Riazuddin and Fayyazuddin, \emph{Phys. Rev.} \textbf{147}, 1071 (1966).

\item \phantomsection\label{ch20-ref-7}
M. Ademollo, G. Veneziano, and S. Weinberg, \emph{Phys. Rev. Lett.} \textbf{22}, 83 (1969).

\item \phantomsection\label{ch20-ref-8}
J. F. Donoghue and E. Golowich, \emph{Phys. Rev.} \textbf{D49}, 1513 (1994).

\item \phantomsection\label{ch20-ref-9}
This result was reported at the 1968 `Rochester' conference at Vienna, and published in E. D. Bloom \emph{et al.}, \emph{Phys. Rev. Lett.} \textbf{23}, 930 (1969); M. Breidenbach \emph{et al.}, \emph{Phys. Rev. Lett.} \textbf{23}, 935 (1969).

\item \phantomsection\label{ch20-ref-10}
J. D. Bjorken, \emph{Phys. Rev.} \textbf{179}, 1547 (1969).

\item \phantomsection\label{ch20-ref-11}
R. P. Feynman, \emph{Phys. Rev. Lett.} \textbf{23}, 1415 (1969); \emph{Photon-Hadron Interactions} (Benjamin, New York, 1972).

\item \phantomsection\label{ch20-ref-12}
C. G. Callan and D. J. Gross, \emph{Phys. Rev. Lett.} \textbf{22}, 156 (1969). The reader should be warned that the symbol $\omega$ as used by Callan and Gross is $2/\omega$ in the notation used here.

\item \phantomsection\label{ch20-ref-13}
D. Gross and S. Treiman, \emph{Phys. Rev.} \textbf{D4}, 1059 (1971).

\item \phantomsection\label{ch20-ref-14}
N. Christ, B. Hasslacher, and A. H. Mueller, \emph{Phys. Rev.}, \textbf{D6}, 3543 (1972).

\item \phantomsection\label{ch20-ref-15}
G. Altarelli and G. Parisi, \emph{Nucl. Phys.} \textbf{B126}, 298 (1977).

\item \phantomsection\label{ch20-ref-16}
H. Georgi and H. D. Politzer, \emph{Phys. Rev.} \textbf{D9}, 416 (1974).

\item \phantomsection\label{ch20-ref-17}
D. J. Gross and F. Wilczek, \emph{Phys. Rev.} \textbf{D9}, 980 (1974).

\item \phantomsection\label{ch20-ref-18}
F. J. Dyson, \emph{Phys. Rev.} \textbf{85}, 631 (1952).

\item \phantomsection\label{ch20-ref-19}
See, e.g., G. N. Hardy, \emph{Divergent Series} (Oxford University Press, Oxford, 1949).

\item \phantomsection\label{ch20-ref-20}
L. N. Lipatov, Leningrad Nuclear Physics Institute report, 1976 (unpublished); \emph{Proceedings of the XVIII International Conference on High Energy Physics at Tbilisi}, 1976.

\item \phantomsection\label{ch20-ref-21}
G. 't Hooft, in \emph{The Whys of Subnuclear Physics -- Proceedings of the 1977 Erice Summer School}, ed. A. Zichichi (Plenum, New York, 1978). For a collection of articles on renormalons and high orders of perturbation theory, see \emph{Large Order Behavior of Perturbation Theory}, eds. J. C. Le Guillou and J. Zinn-Justin (North-Holland, Amsterdam, 1990). For further work on infrared renormalons, see A. H. Mueller, \emph{Phys. Lett. B} \textbf{308}, 355 (1993); \emph{Nucl. Phys.} \textbf{B250}, 327 (1995); A. Duncan and S. Pernice, \emph{Phys. Rev. D} \textbf{51}, 1956 (1995), and articles quoted therein.

\item \phantomsection\label{ch20-ref-22}
B. Lautrup, \emph{Phys. Lett.} \textbf{76B}, 109 (1977).
\end{enumerate}
